\subsection{Энергия}
%%%%%%%%%%%%%%%%%%%%%%%%%%%%%%%%%%%%%%%%%%%%%%%%%%%%
%%%%%%%%%%%%%%%%%%%%%%%%%%%%%%%%%%%%%%%%%%%%%%%%%%%%
%%%%%%%%%%%%%%%%%%%%%%%%%%%%%%%%%%%%%%%%%%%%%%%%%%%%
\z{На неподвижное тело массой $m = 0.5 \,\text{кг}$ начинает действовать постоянная сила $F = 2 \,\text{Н}$. Найти кинетическую энергию, которой будет обладать тело через время $t = 3 \,\text{с}$ после начала действия силы.
}
%
{$E_\text{к} = \frac{F^2 t^2}{2m} = 36 \,\text{Дж}$}

%%%%%%%%%%%%%%%%%%%%%%%%%%%%%%%%%%%%%%%%%%%%%%%%%%%%
\z{Санки массой $m = 2 \,\text{кг}$ и длиной $l = 1 \,\text{м}$ выезжают со льда на асфальт. Коэффициент трения полозьев об асфальт $\mu = 0.5$. Какую работу совершит сила трения к моменту, когда санки полностью окажутся на асфальте?
}
%
{$A_\text{тр} = -\frac{\mu m g l}{2} = -5 \,\text{Дж}$}

%%%%%%%%%%%%%%%%%%%%%%%%%%%%%%%%%%%%%%%%%%%%%%%%%%%%
\z{Пружину, жёсткость которой $k = 200 \,\text{Н/м}$, растянули на одну третью часть её длины, длина пружины в недеформированном состоянии $l_0 = 30 \,\text{см}$. Найти потенциальную энергию пружины.
}
%
{$E_\text{п} = \frac{\alpha^2 k l_0^2}{2} = 1 \,\text{Дж}$}

%%%%%%%%%%%%%%%%%%%%%%%%%%%%%%%%%%%%%%%%%%%%%%%%%%%%
\z{Какую минимальную работу надо совершить, чтобы из колодца глубины $H = 10 \,\text{м}$ поднять на тросе ведро с водой массой $m = 8 \,\text{кг}$? Линейная плотность троса $\mu = 0.4 \,\text{кг/м}$.
}
%
{$A_{\min} = \left(m + \frac{\mu H}{2}\right) g H = 10^3 \,\text{Дж}$}

%%%%%%%%%%%%%%%%%%%%%%%%%%%%%%%%%%%%%%%%%%%%%%%%%%%%
\z{Чтобы удалить гвоздь длиной $10 \,\text{см}$ из бревна, необходимо приложить начальную силу $2 \,\text{кН}$. Гвоздь вытащили из бревна. Какую при этом совершили механическую работу?
}
%
{$A_\text{мех} = \frac{F_0 l}{2} = 100 \,\text{Дж}$}

%%%%%%%%%%%%%%%%%%%%%%%%%%%%%%%%%%%%%%%%%%%%%%%%%%%%
\z{Чтобы растянуть пружину на $3 \,\text{см}$, надо совершить работу $30 \,\text{Дж}$. На сколько растянулась бы пружина, если бы совершили работу $120 \,\text{Дж}$?
}
%
{$x_2 = x_1 \sqrt{\frac{A_2}{A_1}} = 6 \,\text{см}$}

%%%%%%%%%%%%%%%%%%%%%%%%%%%%%%%%%%%%%%%%%%%%%%%%%%%%
\z{Рабочий поднял с помощью подвижного блока груз массой $m = 40 \,\text{кг}$ на высоту $h = 10 \,\text{м}$, прикладывая силу $F = 250 \,\text{Н}$. Определите полезную работу, совершённую работу и КПД механизма.
}
%
{$A_\text{п} = m g h = 4000 \,\text{Дж}, \quad A_c = 2 F h = 5000 \,\text{Дж}, \quad \eta = \frac{m g}{2 F} = 0.8$}

